\documentclass[11pt,a4paper]{moderncv}
\moderncvstyle{banking}
\moderncvcolor{blue}
\usepackage[scale=0.75]{geometry}
\usepackage[utf8]{inputenc}
\usepackage[T1]{fontenc}
\usepackage{hyperref}

% -----------------------------
% Datos personales
% -----------------------------
\name{Matihus}{Molina Larios}
\address{Lima, Perú}{}
\email{matihus.molina@unmsm.edu.pe}
\social[linkedin]{maticus}
\social[github]{Maticus-7}

% -----------------------------
\begin{document} 
	\makecvtitle
	
	% -----------------------------
	% Perfil profesional
	% -----------------------------
	\section{Perfil profesional}
	\cvitem{}{Matemático interdisciplinario apasionado por cimentar las bases teóricas de la Inteligencia Artificial a través de los Sistemas Dinámicos. Combino el rigor analítico puro con la curiosidad computacional para traducir matemáticas abstractas en soluciones aplicables. Mi misión al cursar estudios de posgrado es convertirme en un líder científico que impulse la investigación rigurosa en IA y catalice la autonomía tecnológica de Latinoamérica.}
	
	% -----------------------------
	% Educación
	% -----------------------------
	\section{Educación}
	\cventry{03/2022 -- 12/2026 (estimado)}{Estudiante de Bachillerato en Ciencias Matemáticas}{Universidad Nacional Mayor de San Marcos}{Lima, Perú}{}{Promedio general (CGPA): 15.324/20}
	
	% -----------------------------
	% Idiomas
	% -----------------------------
	\section{Idiomas}
	\cvitem{Español}{Nativo}
	\cvitem{Inglés}{Intermedio B1 (Egresado del Centro Cultural LIDEMPERÚ. En preparación para certificación internacional)}
	\cvitem{Portugués}{Básico A2}
	
	% -----------------------------
	% Proyectos académicos
	% -----------------------------
	\section{Proyectos académicos}
	
	\cventry{}{Investigación Actual Independiente}{Numerical solution for singular equations of Lane-Emden type by ADM and Finite volume method}{}{}{
		\begin{itemize}
			\item Extensión del teorema de existencia y unicidad para el problema de Lane-Emden generalizado no homogéneo mediante un sistema EDO de primer orden $3 \times 3$.
			\item Estimación numérica mediante métodos de volúmenes finitos y contraste con Redes Neuronales.
			\item \href{https://vii.imms.pe/documentos/26PosterMolinaLarios.pdf}{Haga click aquí para ver el póster oficial (Presentado en el VII EICM)}
			\item
			\href{https://maticus-7.github.io/lane-emden/}{Haga click aquí para ver el repositorio en GitHub}
	\end{itemize}}	
	
	\cventry{02/2026 -- 03/2026}{Optimización Numérica y Aprendizaje de Operadores}{SaC$^3$ -- Summer School and Workshop}{}{}{
		\begin{itemize}
			\item Asesor: Leonidas Gkimisis.
			\item Implementé y documenté en Python algoritmos de optimización matemática analizados en el taller, aplicando control de versiones con Git/GitHub para asegurar código reproducible.
			\item \href{https://github.com/Maticus-7/Data-Science}{Haga click aquí para ver el repositorio en GitHub}
	\end{itemize}}
	
	\cventry{01/2025 -- 03/2025}{Proyecto: ``Un primer encuentro con los sistemas dinámicos caóticos''}{PICV-UNMSM}{}{}{
		\begin{itemize}
			\item Docente asesor: Dr. Jorge L. Crisóstomo.
			\item Estudio de la conjugación topológica entre $Q_c$ y el shift de Bernoulli, aplicando conceptos de ecuaciones diferenciales, análisis funcional y topología.
			\item \href{http://bit.ly/3Hj8NvJ}{Haga click aquí para ver el repositorio en GitHub}
	\end{itemize}}
	
	% -----------------------------
	% Eventos académicos e Internacionalización
	% -----------------------------
	\section{Movilidad Internacional y Eventos Académicos}
	
	\cventry{Septiembre 2026}{Heidelberg, Alemania}{13th Heidelberg Laureate Forum }{HLF}{}{Seleccionado mundialmente como \textit{Young Researcher}. Participación en el foro internacional que reúne a los ganadores de la Medalla Fields, Premio Abel y Premio Turing con jóvenes investigadores excepcionales.}
	
	\cventry{Agosto 2026}{Santiago, Chile}{Programa ``San Marquinos por el Mundo''}{Universidad de Chile}{}{Pasantía académica internacional seleccionada por orden de mérito (Top 2 de E.P.). Curso de especialización (40 hrs): \textit{Desafíos Globales, Transformación Digital e Innovación}.}
	
	\cventry{Diciembre 2025}{Huaraz, Perú}{VII Encuentro Internacional de Ciencias Matemáticas}{EICM}{}{Ponente. Presentación de póster de investigación sobre soluciones numéricas para la ecuación de Lane-Emden.}
	
	% -----------------------------
	% Habilidades Técnicas
	% -----------------------------
	\section{Habilidades Técnicas}
	\cvitem{Lenguajes y Entornos}{Python, R, MATLAB, LEAN4, Git, \LaTeX, Power BI.}
	\cvitem{Librerías}{NumPy, SciPy, Pandas, Scikit-Learn, PyTorch, Matplotlib.}
	
	% ------------------------------
	% Logros
	% ------------------------------
	\section{Premios y Distinciones}
	\cvitem{2025}{Beca Permanencia de Estudios Nacionales, otorgada por PRONABEC por alto rendimiento académico.}
	\cvitem{2025}{Reconocimiento de Excelencia Académica otorgado por la Facultad de Ciencias Matemáticas de la UNMSM.}  
	
	% -----------------------------
	% Actividades extracurriculares
	% -----------------------------
	\section{Liderazgo y Actividades Extracurriculares}
	\cvitem{Actual}{Director del área de Proyectos e Investigación del IEEE Computational Intelligence Society (UNMSM) y Traine Researcher del grupo BILDHEIT.}
	\cvitem{Pasado}{Exmiembro del equipo universitario de judo (UNMSM) y de la agrupación artística SALSUNI.}
	
	% -----------------------------
	% Certificados
	% -----------------------------
	\section{Certificados y Portafolio} 
	\cvitem{}{ \href{https://maticus-7.github.io/Certificates/}{Los certificados de cursos, Lattes y logros académicos están disponibles en mi repositorio digital (Click aquí)}.}
	
\end{document}